%  07_discussion.tex
\section{Discussion}
\label{sec:discussion}

\subsection{Why pair-based models win at small identity counts}

The consistent advantage of the Siamese network over ArcFace and AdaFace across all
experimental conditions is not an architectural accident but a structural consequence of
the training objective.

Angular-margin losses optimise a \emph{closed-set classification} problem: at training
time the model learns to separate the specific identities it has seen, placing them into
distinct cones on the unit hypersphere.  The hope at test time is that the manifold
structure learned from a few training identities will generalise---that unseen test
identities will land in distinguishable regions of the same embedding space.  With only
6 training identities and 512 embedding dimensions, the hypersphere has vastly more
capacity than necessary for the training task, so the learned geometry is underconstrained
and effectively arbitrary with respect to novel bats.

The Siamese network optimises a \emph{pairwise verification} objective directly: for
every training batch the model is asked whether two images show the same bat.  It never
assigns identities to fixed slots and therefore does not depend on the angular geometry
generalising to unseen individuals.  This alignment between training objective and
deployment task is the structural reason the pair-based model wins.

We expect this ordering to reverse as identity count grows: with many training identities
the angular manifold becomes well-constrained and embedding-based models typically
dominate~\cite{deng2019arcface, kim2022adaface}.

\subsection{Evidence that the current identity count is insufficient}
\label{sec:identity_wall}

Four independent signals in the experimental data converge on the same conclusion: the
dataset is currently too small for angular-margin losses to learn transferable geometry.

\paragraph{Train/test performance divergence.}
The sharpest signal is the gap between training accuracy and test ROC-AUC for the
embedding models:

\begin{table}[h]
\centering
\caption{Training accuracy vs.\ test ROC-AUC for embedding models (rousettus).
Near-perfect closed-set training accuracy coincides with near-chance test discrimination.}
\label{tab:overfit}
\setlength{\tabcolsep}{8pt}
\begin{tabular}{lccc}
\toprule
Model & Train acc & Test ROC-AUC & Gap \\
\midrule
ArcFace (tuned) & 0.988 & 0.541 & $-0.447$ \\
AdaFace (tuned) & 1.000 & 0.507 & $-0.493$ \\
Siamese         & 1.000 & 0.809 & $-0.191$ \\
\bottomrule
\end{tabular}
\end{table}

ArcFace and AdaFace perfectly separate the 6 training identities yet produce test
ROC-AUC barely above chance.  The Siamese network, while also fully fitting the training
set, retains a much smaller gap because the pairwise objective directly optimises the
verification criterion being evaluated.

\paragraph{Top-1 identification at chance level.}
With 3 unseen test identities, random guessing yields a Top-1 accuracy of $0.33$.
ArcFace achieves $0.344$ and AdaFace $0.286$---both within noise of chance.
The embedding space has no useful structure for assigning novel identities;
it has only memorised the angular positions of the 6 training cones.

\paragraph{Optimization instability.}
The validation ROC-AUC curves for ArcFace and AdaFace (Figure~\ref{fig:learning_curves})
show high epoch-to-epoch variance throughout training, never stabilising, in contrast
to the Siamese network which converges smoothly by epoch 4.  This instability is a
symptom of an under-constrained optimisation landscape: with too few classes, the
angular loss surface has many shallow plateaus and the model cannot find a configuration
that generalises.  The sensitivity to batch ordering observed before the DataLoader seed
fix (repro-1 vs.\ repro-2: ROC-AUC $0.528$ vs.\ $0.563$ from identical hyperparameters)
further illustrates this fragility.

\paragraph{Cross-species consistency rules out species-specific explanations.}
The same failure pattern appears independently for both species:

\begin{table}[h]
\centering
\caption{ArcFace test ROC-AUC across species (random background, tuned hyperparameters).
The near-chance performance and the Siamese advantage are consistent across both
species, pointing to identity count as the common cause.}
\label{tab:crossspecies}
\setlength{\tabcolsep}{7pt}
\begin{tabular}{lcccc}
\toprule
Species & Train ids & Siamese & ArcFace & Gap \\
\midrule
rousettus & 6 & 0.809 & 0.541 & $+0.268$ \\
mauritius & 5 & 0.803 & 0.526 & $+0.277$ \\
\bottomrule
\end{tabular}
\end{table}

Had the failure been caused by species-specific factors---fur texture,
pose distribution, image quality---the gap would vary between species.
Instead it is virtually identical ($0.268$ vs.\ $0.277$), which isolates
identity count as the governing variable.  This cross-species replication
is the strongest experimental evidence that the current scale is the bottleneck
rather than any property of the images or animals themselves.

\subsection{Background effects}

The green-screen and original-background conditions outperform random-background
compositing for both species on the Siamese network ($0.925/0.922$ vs.\ $0.803$ on
mauritius; Table~\ref{tab:mauritius}).  Two explanations are plausible: (1) random
texture compositing introduces visual clutter that competes with facial features in the
$105\!\times\!105$ input field; (2) consistent backgrounds reduce domain shift within
the test split.

The embedding models show a stronger relative improvement under green-screen ($0.761$
for ArcFace green vs.\ $0.526$ for ArcFace random on mauritius), suggesting that the
angular-margin geometry is more sensitive to background variation than the pairwise
objective.

\subsection{The leaky-vs-honest gap}

The 24-point F\textsubscript{1} gap between the random-pair protocol and the
identity-disjoint protocol (Table~\ref{tab:leakage}) is larger than typical published
figures for animal re-identification.  Practitioners are encouraged to audit their
evaluation protocols: if the same individual appears in both train and test, reported
numbers are not measuring deployed-scenario performance.

\subsection{Limitations}

\paragraph{Single-seed evaluation.}
All reported numbers come from single training runs.  The standard deviation of test
metrics across seeds is unknown and likely non-trivial for the embedding models given
the instability documented in Section~\ref{sec:identity_wall}.

\paragraph{No external validation set.}
The test identities were drawn from the same colony as the training identities.
Generalisation to bats from different colonies or imaging setups is untested.

\paragraph{No ablation on backbone pretraining.}
All embedding models used ImageNet-pretrained ResNet-50 weights.  Whether
domain-adaptive pretraining on a larger bat-face corpus would lower the identity-count
threshold is an open question.
